% !Mode:: "TeX:UTF-8"
\begin{acknowledgements}
近三年硕士阶段的学习与科研工作已接近尾声，回望这段历程，既有攻克难题时的欣喜，也有反复推敲与沉潜钻研的不易。值此论文完成之际，谨向所有在学位论文的研究与写作过程中给予我支持与帮助的师长、同门及家人致以最诚挚的感谢。

首先，谨向我的导师王莘教授致以崇高的敬意和衷心的感谢。王教授学识渊博、治学严谨，在本课题的立项、研究方法选择与实施、数据分析以及论文撰写等各个环节给予了悉心指导与严格把关。三年来，王莘教授不仅传授了宝贵的科研方法与学术规范，更以身体力行的治学态度深刻影响了我对待研究的方式，这将是我终身受用的财富。特别感谢陆晔师兄在课题研究方向、实验设计与改进思路等具体环节提供的耐心帮助和宝贵建议，很多关键问题在与师兄的深入讨论中得以厘清和解决。感谢实验室各位同门在日常科研中的学术交流、互相鼓励与共同成长，课题组温暖而严谨的氛围是我坚持下去的重要动力之一。

感谢哈尔滨工业大学计算学部及网络空间安全学院在三年间给予的悉心培养与充分支持。学校为研究生提供了先进的科研平台与充足的计算资源，校内高性能计算环境为本研究中大量深度学习实验的顺利开展提供了坚实保障。课程学习阶段，各位任课教师严谨的教学态度与深厚的学术功底为我打下了坚实的理论基础。哈尔滨工业大学"规格严格，功夫到家"的校训精神如同一条红线贯穿于三年学习始终，督促我在面对研究难点时不敷衍、不妥协，以高标准要求自己的每一项工作。这段求学时光，是我人生中弥足珍贵的经历，学校浓厚的学术氛围与严谨求实的科研文化将深远地影响我今后的工作与研究。感谢参与本论文审阅与答辩的各位专家和老师，感谢你们在百忙之中对论文进行仔细审阅，所提出的宝贵意见和建议使论文质量得到了显著提升。

最后，谨向我的家人表示最衷心的感谢。求学三年，父母始终是我最坚实的后盾，他们在生活上无微不至的关怀与默默的付出，让我得以全身心投入学业与科研，无后顾之忧。感谢所有亲友在精神上给予的理解、包容与鼓励，每一份温暖的问候都是我在艰难时刻坚持前行的力量。对所有曾给予我帮助与支持的人，谨致以最诚挚的谢意。
\end{acknowledgements}